\subsection{On-Off Control Performance}
\label{subsec:on-off-control-performance}

The circuit was initially tested using a basic on-off control strategy to
maintain a target temperature of \(55^\circ\)C. As shown in
Figure~\ref{fig:baseline}, the system's temperature increased steadily over
time and eventually leveled off near the target temperature.

\begin{figure}[H]
    \centering
    \includegraphics[width=0.8\textwidth]{../out/baseline}
    \caption{
        Temperature vs. time for the on-off control circuit regulating an
        aluminum block to a target of $55^\circ$C. The system shows an
        oscillating temperature curve, with local minima just below the set
        temperature, demonstrating the basic performance of simple on-off
        control.
    }
    \label{fig:baseline}
\end{figure}

When hysteresis control was introduced, with a 3°C range around the target
temperature, the system exhibited a pattern of temperature fluctuations, as
shown in Figure~\ref{fig:hysteresis}.

\begin{figure}[H]
    \centering
    \includegraphics[width=0.8\textwidth]{../out/simple_hysteresis}
    \caption{
        Temperature vs. time for the on-off circuit with a 3$^\circ$C
        hysteresis band and a $55^\circ$C target. The plot shows
        characteristic oscillations as the heater cycles on and off,
        with sharp increases during heating phases and less pronounced drops
        during cooling.
    }
    \label{fig:hysteresis}
\end{figure}

The temperature increased steadily during heating but exceeded the upper limit
of the range before cooling down rapidly.
The temperature then dropped below the lower limit of the range before the
heater turned back on.
This uneven heating and cooling pattern suggests that the aluminum block loses
heat slower than it absorbs heat from the heater.
These results highlight the limited accuracy of simple on-off control methods
for maintaining a stable temperature in systems like this.
These measurements provided a baseline for comparing the performance of more
advanced control methods.

\subsection{Proportional Control Implementation}
\label{subsec:proportional-control-implementation2}

Proportional control was implemented to improve temperature regulation.
This method adjusts the heater's power based on how far the current
temperature is from the target.
Figure~\ref{fig:proportional_3_0.1} shows the system's behavior with a
proportionality factor (\(K_p\)) of 0.3, an offset of 0.1, and a temperature
range of 3°C. This approach reduced the size of temperature overshoots and
made the system more stable, although the offset value was not ideal for
keeping the temperature close to the target.

\begin{figure}[H]
    \centering
    \includegraphics[width=0.8\textwidth]{../out/proportional_0.1_offset}
    \caption{
        Temperature vs. time for proportional control with a 3$^\circ$C band,
        offset of 0.1, and $K_p = 0.3$ for an aluminum block at $55^\circ$C.
        The system exhibits reduced overshoot and improved stability compared
        to on-off control, but the offset is not fully optimized.
    }
    \label{fig:proportional_3_0.1}
\end{figure}

A further adjustment to the offset, increasing it to 0.165, improved the
performance, as shown in Figure~\ref{fig:proportional_3_0.165}.

\begin{figure}[H]
    \centering
    \includegraphics[width=0.8\textwidth]{../out/proportional_0.165_offset}
    \caption{
        Temperature vs. time for proportional control with a 3$^\circ$C band,
        offset of 0.165, and $K_p = 0.3$ for an aluminum block at $55^\circ$C.
        Improved offset tuning results in better temperature regulation and
        smaller fluctuations.
    }
    \label{fig:proportional_3_0.165}
\end{figure}

The system was able to stay within the target range for longer periods, and
the temperature fluctuations became smaller over time.
This behavior is typical of systems using proportional control, where the
heater's power is adjusted based on the difference between the current
temperature and the target, allowing for smoother changes compared to simple
on-off control.

The temperature range was then reduced to 1°C while keeping the same optimized
settings (\(K_p = 0.3\), offset = 0.165).
This adjustment allowed the system to regulate the temperature more precisely,
as shown in Figure~\ref{fig:proportional_1_0.165}.

\begin{figure}[H]
    \centering
    \includegraphics[width=0.8\textwidth]{../out/proportional_0.165_offset_1_band}
    \caption{
        Temperature vs. time for proportional control with a 1$^\circ$C band,
        offset of 0.165, and $K_p = 0.3$ for an aluminum block at $55^\circ$C.
        Narrowing the control band allows for more precise temperature
        regulation, though a small error remains.
    }
    \label{fig:proportional_1_0.165}
\end{figure}

However, the system still exhibited a small but consistent error, where the
temperature was slightly higher or lower than the target.
This occurred because proportional control requires a small difference between
the actual and target temperatures to maintain the heater's power and balance
heat loss to the environment.

\subsection{Proportional-Integral Control}
\label{subsec:proportional-integral-control}

To address the consistent temperature error observed with proportional
control, an integral control component was added.
Figures~\ref{fig:integral_1_0.165},
and~\ref{fig:live_plot_20251027_153510_integral}, show the performance of the
proportional-integral (PI) controller with \(K_p = 0.3\), \(K_i = 0.00001\),
an offset of 0.165, and a 1°C range, applied to the aluminum block at a target
temperature of \(55^\circ\)C. The integral control component works by
gradually correcting any long-term differences between the current temperature
and the target.

\begin{figure}[H]
    \centering
    \includegraphics[width=0.8\textwidth]{../out/integral_0.165_offset_1_band}
    \caption{
        Temperature vs. time for PI control with a 1$^\circ$C band, offset of
        0.165, $K_p = 0.3$, and $K_i = 0.00001$ for an aluminum block at
        $55^\circ$C. The addition of integral action reduces steady-state
        error and improves long-term stability.
    }
    \label{fig:integral_1_0.165}
\end{figure}

\begin{figure}[H]
    \centering
    \includegraphics[width=0.8\textwidth]{../out/live_plot_20251027_153510_integral}
    \caption{
        Temperature vs. time for PI control (Aluminum):
        $T_\text{set}=55^\circ$C, band=1, $K_p=0.3$, $K_i=0.00001$,
        offset=0.165. The plot demonstrates the system's ability to maintain
        temperature near the setpoint over time.
    }
    \label{fig:live_plot_20251027_153510_integral}
\end{figure}

This adjustment allowed the system to reduce the consistent error and
stabilize the temperature closer to the desired value.
Over time, the size of the temperature fluctuations decreased until the system
maintained a more stable performance near the target.

At a higher target temperature of \(65^\circ\)C, the PI controller with an
increased integral gain (\(K_i = 0.001\)) and an offset of 0.25 showed
effective performance, as illustrated in
Figure~\ref{fig:live_plot_20251029_163401}.

\begin{figure}[H]
    \centering
    \includegraphics[width=0.8\textwidth]{../out/live_plot_20251029_163401}
    \caption{
        Temperature vs. time for PI control at a higher setpoint ($65^\circ$C)
        with $K_p=0.3$, $K_i=0.001$, $K_d=0.0$, and offset=0.25. The increased
        integral gain enables faster correction of errors, demonstrating
        adaptability to different target temperatures.
    }
    \label{fig:live_plot_20251029_163401}
\end{figure}

The higher integral gain allowed faster correction of temperature errors,
demonstrating that the system could adapt effectively to different target
temperatures.

\subsection{Full PID Control and Derivative Action}
\label{subsec:full-pid-control-and-derivative-action}

The final step in the control strategy involved adding a derivative control
component, completing the proportional-integral-derivative (PID) system.
Derivative control helps the system respond more quickly to sudden changes in
temperature by predicting how the temperature is likely to change.
Figure~\ref{fig:live_plot_20251027_171632_derivative} shows the system's
behavior with the derivative gain set to zero (\(K_d = 0.0\)), effectively
functioning as a PI controller for comparison.

\begin{figure}[H]
    \centering
    \includegraphics[width=0.8\textwidth]{../out/live_plot_20251027_171632_derivative}
    \caption{
        Temperature vs. time for PID control (Aluminum):
        $T_\text{set}=55^\circ$C, band=1, $K_p=0.3$, $K_i=0.00001$, $K_d=0.0$,
        offset=0.165. With $K_d=0$, the system acts as a PI controller,
        providing a baseline for comparison with full PID control.
    }
    \label{fig:live_plot_20251027_171632_derivative}
\end{figure}

When derivative control was activated (\(K_d = 0.00001\)), the system
demonstrated a much faster response to sudden temperature changes, as shown in
Figure~\ref{fig:live_plot_20251028_162354}.

\begin{figure}[H]
    \centering
    \includegraphics[width=0.8\textwidth]{../out/live_plot_20251028_162354}
    \caption{
        Temperature vs. time for PID control (Aluminum):
        $T_\text{set}=55^\circ$C, band=1, $K_p=0.3$, $K_i=0.00001$,
        $K_d=0.00001$, offset=0.165. Activating derivative control results in
        a faster response to sudden temperature changes, maintaining a much
        tighter oscilation around the setpoint.
    }
    \label{fig:live_plot_20251028_162354}
\end{figure}

For instance, during testing, the temperature sensor was removed from the
aluminum, causing it to detect a rapid drop in temperature.
The derivative control responded by immediately increasing the heater power to
100\%, highlighting its ability to handle unexpected changes.
This response occurred with \(K_d = 0.001\), showing that even small
derivative gains have a significant effect on how the system reacts to sudden
events.

However, the derivative control also introduced a drawback: it amplified
small fluctuations in the temperature readings caused by sensor noise.
This resulted in slightly more variability in the heater's power during normal
operation, which partially reduced the benefits of the faster response to
temperature changes.

The PID controller was also tested on a brass block at \(65^\circ\)C, using
\(K_p = 0.3\), \(K_i = 0.001\), \(K_d = 0.001\), and an offset of 0.25, as
shown in Figure~\ref{fig:live_plot_20251029_170448}.

\begin{figure}[H]
    \centering
    \includegraphics[width=0.8\textwidth]{../out/live_plot_20251029_170448}
    \caption{
        Temperature vs. time for PID control (Brass):
        $T_\text{set}=65^\circ$C, band=1, $K_p=0.3$, $K_i=0.001$, $K_d=0.001$,
        offset=0.25. The PID controller successfully regulates the temperature
        of a brass block, demonstrating adaptability to different materials.
    }
    \label{fig:live_plot_20251029_170448}
\end{figure}

The PID controller successfully regulated the temperature of the brass block,
demonstrating its ability to adapt to materials with different thermal properties.

\subsection{Specific Heat Capacity Measurements}
\label{subsec:specific-heat-capacity-measurements}

The specific heat capacity determinations yielded results with substantial
deviations from accepted literature values.
For the aluminum block (\(24.64 \pm 0.005\) g), the calculated specific heat
capacity was \(1.8 \, \frac{\text{J}}{\text{g}\cdot{}^\circ\text{C}}\),
representing a 100\% error relative to the accepted value of
\(0.9 \, \frac{\text{J}}{\text{g}\cdot{}^\circ\text{C}}\)~\cite{specific_heat}.
Similarly, the brass block (\(72.35 \pm 0.005\) g) yielded
\(0.8 \, \frac{\text{J}}{\text{g}\cdot{}^\circ\text{C}}\) compared to the
literature value of \(0.4 \, \frac{\text{J}}{\text{g}\cdot{}^\circ\text{C}}\)
~\cite{specific_heat}, again corresponding to a 100\% error.

While it was noted that aluminum measurements typically exceed
theoretical values in this experimental setup, the magnitude of the
discrepancies suggests systematic error rather than random uncertainty.
Potential error sources include: (1) inaccuracies in PWM-to-power conversion
calculations, (2) errors in MOSFET switching efficiency modeling, (3)
miscalibration of the energy delivery measurement, (4) unaccounted heat losses
to the environment during the heating interval, or (5) noise amplification in
the data acquisition chain affecting temperature change rate calculations.

Experimental limitations prevented comprehensive error analysis through
extended data collection.
Hardware failures during the experiment, including damaged components
requiring replacement, significantly compressed the available testing time.
A more extensive measurement campaign spanning a wider temperature range would
have facilitated identification of whether errors exhibited systematic trends
or random scatter, thereby isolating the dominant uncertainty source.
The consistent overestimation for both materials suggests a systematic
calibration error in the power delivery measurement or energy calculation
algorithm rather than material-specific measurement artifacts.
